% Generated draft: requires mathematical and source review.
\documentclass[11pt]{article}
\usepackage[margin=25mm]{geometry}
\usepackage{fontspec}
\setmainfont{Latin Modern Roman}
\usepackage{amsmath,amssymb,mathtools,unicode-math}
\setmathfont{Latin Modern Math}
\usepackage{enumitem,needspace,xurl,hyperref}
\setlength{\parindent}{0pt}
\setlength{\parskip}{5pt}
\setlength{\emergencystretch}{3em}
\newcommand{\N}{\mathbb{N}}
\newcommand{\Z}{\mathbb{Z}}
\newcommand{\Q}{\mathbb{Q}}
\newcommand{\R}{\mathbb{R}}
\newcommand{\C}{\mathbb{C}}
\newcommand{\F}{\mathbb{F}}
\newcommand{\A}{\mathbb{A}}
\newcommand{\PP}{\mathbb{P}}
\newcommand{\E}{\mathbb{E}}
\newcommand{\eps}{\varepsilon}
\newcommand{\dd}{\,\mathrm{d}}
\begin{document}
\section*{UnsolvedMath ID: 30001769}
\textbf{Source problem code:} catalogue entry
\subsection{Definitions and mathematical statement}
To explain the terminology, fix a positive integer \(n\), a prime \(p\), and a field \(k\) of characteristic \(p\), and let \(G=\mathfrak S_n\). These are parameters for the definitions; the supplied record does not specify the intended field or the range of these parameters. All modules in this explanatory framework are finite-dimensional left modules.

An indecomposable \(kG\)-module is a nonzero module that is not a direct sum of two nonzero submodules. A \(kG\)-module \(M\) is relatively projective with respect to a subgroup \(Q\leq G\) if it is isomorphic to a direct summand of \(\operatorname{Ind}_Q^G\operatorname{Res}_Q^G M\). A vertex of an indecomposable \(kG\)-module \(M\) is a \(p\)-subgroup \(Q\leq G\) minimal under inclusion among those for which \(M\) is relatively projective. Vertices of a fixed indecomposable module are conjugate in \(G\), so their determination is understood up to \(G\)-conjugacy.

For a partition \(\lambda\vdash n\), the Specht module \(S_k^\lambda\) is the submodule of the permutation module on \(\lambda\)-tabloids spanned by the polytabloids
\[
 e_t=\sum_{\sigma\in C_t}\operatorname{sgn}(\sigma)\,\{\sigma t\},
\]
where \(t\) ranges over tableaux of shape \(\lambda\), \(C_t\) is the subgroup permuting entries within each column of \(t\), and \(\{t\}\) identifies tableaux differing only by permutations within rows. The coefficients are interpreted in \(k\).

\textbf{Source classification request.} Determine the vertices of naturally occurring modular representations of symmetric groups, particularly Specht modules, beyond the known special cases.

This is a research direction rather than a fully specified conjecture. For a chosen indecomposable module in an intended family, the mathematical target is the conjugacy class of its vertices. For an indecomposable Specht module this applies directly to \(S_k^\lambda\). If a Specht module is decomposable, applying the definition above requires choosing its indecomposable summands; interpreting the request as asking for the vertices of all such summands is a possible clarification, not an explicit requirement in the source. The record does not specify the family meant by ``naturally occurring,'' the intended parameters, the treatment of decomposable modules, or which known special cases are excluded. Consequently, it does not determine a unique classification problem with complete quantifiers.
\subsection{Short English statement}
The supplied problem asks for a classification of vertices, which describe the minimal \(p\)-subgroups relative to which an indecomposable modular representation is projective. Specht modules are the main named family. The source leaves the intended scope and several hypotheses unspecified, so this draft records the classification direction and its unresolved qualifications.
\subsection{Sources}
\begin{itemize}
\item Local registry: \texttt{data/unsolvedmath/problems.json}, permanent UnsolvedMath ID \(30001769\), source problem code \texttt{catalogue entry}, title ``Vertices of Modular Symmetric-Group Representations.'' The supplied record lists the proposed year as 2011 and names no proposer.
\item Supplied report citation: \emph{Mini-Workshop: Modular Representations of Symmetric Groups and Related Objects} (2011), \emph{Oberwolfach Reports}, pp.\ 1177--1205, DOI 10.4171/owr/2011/21. Original OWR report: \url{https://doi.org/10.4171/owr/2011/21}.
\item Supplied literature source, labelled ``Wildon vertices of Specht modules'': \url{https://arxiv.org/abs/0907.0811}. The supplied record gives no further bibliographic details.
\item Literature assessment date: \texttt{2026-08-21}. This is supplied dated curation, not independently verified. It labels the problem open, describes the extraction as an introductory fragment with no determinate conjecture, and reports that no complete classification is established. Its accompanying note requests independent verification of time-sensitive status and cited claims.
\item The registry also labels the statement ``exact'' and reports that manual review found it self-contained and mathematically coherent. These are supplied verification claims; their tension with the fragment assessment remains unresolved. Neither linked source was independently consulted for this draft.
\end{itemize}
\end{document}
